\documentclass[letterpaper,11pt]{article}
\usepackage[T1]{fontenc}
\usepackage[utf8]{inputenc}
\IfFileExists{spanish.ldf}{\usepackage[spanish,es-tabla]{babel}}{\usepackage{babel}\babelprovide[import,main]{spanish}}
\usepackage{lmodern,graphicx,booktabs,tabularx,fvextra}
\usepackage[margin=2.5cm]{geometry}
\usepackage[hidelinks,hypertexnames=false]{hyperref}
\setlength{\parindent}{0pt}
\setlength{\parskip}{0.5em}
\newcommand{\bloque}[3][\linewidth]{%
 \IfFileExists{../evidencias/#2}{%
   \begin{center}
   \begin{minipage}{\linewidth}\centering
   \includegraphics[width=#1,height=0.36\textheight,keepaspectratio]{../evidencias/#2}\par\smallskip
   {\small #3\par}
   \end{minipage}
   \end{center}%
 }{\textit{Archivo de evidencia: \nolinkurl{#2}.}\par}%
}
\begin{document}
\begin{titlepage}
\centering
\vspace*{1.6cm}
{\large Universidad Diego Portales\par}
{Criptografía y Seguridad en Redes\par}
\vspace{2.2cm}
{\Huge\bfseries Informe Laboratorio 4\par}
\vspace{0.4cm}
{\Large SSH y análisis de tráfico cifrado\par}
\vspace{1.2cm}
{\large Francisco Antonio Tapia Canales\par}
\href{mailto:francisco.tapia1@mail.udp.cl}{francisco.tapia1@mail.udp.cl}\par
\vspace{0.5cm}
Sección 4\par
Profesor Brayan Espina\par
\vfill
Laboratorio 4 · 2026\par
\vspace{0.9cm}
\end{titlepage}
\tableofcontents
\clearpage

\section{Descripción de actividades}
La práctica se realizó en una máquina virtual con Ubuntu. Primero se configuró un servidor SSH en Docker y se observó en Wireshark qué información quedaba visible al iniciar la conexión. Luego se analizó la captura entregada por el curso. Con el par de claves RSA de prueba se recuperó el PIN y se derivó la clave que permitió verificar y leer las transacciones.

\section{Desarrollo}

\subsection{Servidor SSH en Docker}
Se construyó la imagen \texttt{udp-lab4-ssh:local} con el identificador \texttt{a9f78de2a076}. El \texttt{Dockerfile} usa Debian Bookworm e instala OpenSSH. Al iniciar, el contenedor generó una clave de host Ed25519 y dejó \texttt{sshd} en primer plano. El puerto 22 del contenedor quedó disponible en \texttt{127.0.0.1:2222} de Ubuntu.

En \texttt{ssh/sshd\_config} se habilitó la autenticación por clave pública para el usuario de prueba. El archivo deshabilitó la autenticación por contraseña y el acceso de \texttt{root}. La salida de \texttt{sshd -T} dentro del contenedor confirmó la siguiente configuración:
\begin{Verbatim}[fontsize=\small]
passwordauthentication no
ciphers aes256-ctr
macs hmac-sha2-256
kexalgorithms curve25519-sha256
hostkeyalgorithms ssh-ed25519
authenticationmethods publickey
\end{Verbatim}

\begin{tabularx}{\linewidth}{@{}lX@{}}
\toprule
\textbf{Algoritmo} & \textbf{Función en la conexión}\tabularnewline
\midrule
\texttt{curve25519-sha256} & Permite acordar un secreto compartido a partir del cual se derivan las claves de sesión.\tabularnewline
\texttt{ssh-ed25519} & Permite que el cliente verifique la firma del servidor durante el intercambio de claves.\tabularnewline
\texttt{aes256-ctr} & Cifra el tráfico con AES de 256 bits en modo contador.\tabularnewline
\texttt{hmac-sha2-256} & Verifica la integridad y autenticidad de los paquetes de la sesión.\tabularnewline
\bottomrule
\end{tabularx}

\subsection{Conexión de prueba}
Se inició la captura de Wireshark en la interfaz \texttt{lo} antes de abrir la conexión. El cliente se autenticó por clave pública y ejecutó un comando remoto. La salida \texttt{LAB4\_SSH\_OK} y el usuario \texttt{laboratorio} confirmaron el acceso. El registro del cliente mostró los algoritmos que se usaron en ambos sentidos:
\begin{Verbatim}[fontsize=\footnotesize,breaklines]
debug1: kex: algorithm: curve25519-sha256
debug1: kex: host key algorithm: ssh-ed25519
debug1: kex: server->client cipher: aes256-ctr MAC: hmac-sha2-256 compression: none
debug1: kex: client->server cipher: aes256-ctr MAC: hmac-sha2-256 compression: none
Authenticated to 127.0.0.1 ([127.0.0.1]:2222) using "publickey".
\end{Verbatim}

\clearpage
\subsection{Tráfico observado en Wireshark}
Las cabeceras TCP mostraron la dirección \texttt{127.0.0.1} en ambos extremos. El cliente usó el puerto \texttt{57062} y el servidor, el \texttt{2222}. Antes del cifrado se leyeron las versiones de OpenSSH: \texttt{10.2p1} en el cliente y \texttt{9.2p1} en el servidor (Figura~1).

\bloque[0.94\linewidth]{ssh-version.png}{Figura 1. Cadenas de identificación SSH en las tramas 132 y 134, con la dirección y los puertos de la conexión.}

Los mensajes \texttt{KEXINIT} mostraron los algoritmos ofrecidos por cada extremo. En la trama 137 se observó la propuesta del servidor (Figura~2). La selección se comprobó por separado con el registro del cliente.

\bloque[0.94\linewidth]{ssh-kexinit.png}{Figura 2. Algoritmos ofrecidos por el servidor en la trama 137.}

\clearpage
Los mensajes \texttt{NEWKEYS} marcaron el paso a las claves negociadas. En la trama 142 apareció un paquete cifrado del cliente (Figura~3). A partir de ese punto, no se pudo leer el contenido de aplicación en la captura. Los puertos, los tamaños de los paquetes y los tiempos de envío siguieron visibles.

\bloque[0.98\linewidth]{ssh-newkeys.png}{Figura 3. Cambio de claves en las tramas 139 y 140 y paquete cifrado del cliente en la trama 142.}

Se exportó el flujo SSH y se comprobó con TShark que la captura resultante contenía 40 tramas de una sola conexión TCP. La exportación cambió los números de trama, como se muestra a continuación.

\begin{tabularx}{\linewidth}{@{}Xcc@{}}
\toprule
\textbf{Evento} & \textbf{Vista original} & \textbf{Archivo exportado}\tabularnewline
\midrule
Identificación del cliente y del servidor & 132 y 134 & 4 y 6\tabularnewline
KEXINIT del cliente y del servidor & 136 y 137 & 8 y 9\tabularnewline
NEWKEYS del servidor y del cliente & 139 y 140 & 11 y 12\tabularnewline
Primer paquete cifrado del cliente & 142 & 14\tabularnewline
\bottomrule
\end{tabularx}

\clearpage
\subsection{Registros RSA y AES-GCM}
En la primera trama de la captura del curso se encontró la etiqueta \texttt{RSA\_PIN|}, seguida por una sal de 16 bytes y el PIN cifrado. Las tramas 2 a 41 contenían los registros \texttt{AESGCM|}. Cada uno llevaba un nonce de 12 bytes y una etiqueta de autenticación de 16 bytes. Se comprobaron 40 nonces distintos.

\begin{tabularx}{\linewidth}{@{}clcX@{}}
\toprule
\textbf{Trama} & \textbf{Registro} & \textbf{Carga TCP} & \textbf{Campos en orden (bytes)}\tabularnewline
\midrule
1 & \texttt{RSA\_PIN|} & 280 & Etiqueta (8), sal (16) y PIN cifrado (256).\tabularnewline
2 & \texttt{AESGCM|} & 119 & Etiqueta (7), nonce (12), texto cifrado (84) y etiqueta GCM (16).\tabularnewline
\bottomrule
\end{tabularx}

\subsection{Recuperación del PIN y derivación de la clave}
Se comprobó que ambas claves de prueba correspondían al mismo par RSA de 2048 bits. La clave privada permitió recuperar el PIN con RSA-OAEP, usando SHA-256 y MGF1 con SHA-256. Se conservaron sus cuatro dígitos como bytes ASCII para que un cero inicial, si lo hubiera, no se perdiera.

Con el PIN y la sal extraída de la trama 1 se aplicó Scrypt. Se usaron \(N=2^{14}\), \(r=8\) y \(p=1\), con una salida de 32 bytes para AES-256. La sal hace que un mismo PIN produzca claves diferentes cuando se cambia ese valor. Por eso fue necesario conservar la sal, los parámetros y la longitud de salida para reproducir la clave. En Scrypt, \(N\) fija el costo de trabajo, \(r\) controla el tamaño del bloque y \(p\) determina el número de bloques que se procesan de forma independiente.

\subsection{Descifrado de las transacciones}
El script \texttt{scripts/descifrar.py} separó las cargas TCP y verificó la etiqueta GCM de cada registro con su propio nonce. La prueba recuperó 40 mensajes y las 40 etiquetas se validaron. Se conservaron el identificador, el remitente, el destinatario y el monto tal como aparecían en los mensajes. El resumen de la ejecución y el inventario de tramas están en \texttt{evidencias/}.
\begin{Verbatim}[fontsize=\small]
RSA: 2048 bits; par de claves verificado.
RSA_PIN: trama 1; PIN de 4 dígitos recuperado.
Scrypt: N=16384, r=8, p=1; sal de 16 bytes; salida de 32 bytes.
AESGCM: 40 mensajes autenticados; 40 nonces distintos.
\end{Verbatim}


\section*{Conclusiones y comentarios}
\addcontentsline{toc}{section}{Conclusiones y comentarios}
En la prueba SSH se pudieron leer las versiones y las propuestas de algoritmos antes del cambio de claves. Después de ese paso, la captura mantuvo visibles los datos de conexión, pero dejó de mostrar el contenido de aplicación. En la segunda actividad, la clave RSA entregada para la práctica permitió recuperar el PIN y, con él, volver a derivar la clave AES con la que se autenticaron las transacciones.

Las claves RSA no permitirían recuperar el secreto de un esquema poscuántico independiente. Se requeriría la clave privada del esquema que se hubiese usado.

\clearpage
\section*{Transacciones recuperadas}
\addcontentsline{toc}{section}{Transacciones recuperadas}
\IfFileExists{../.local/descifrado/transacciones.txt}{%
 \VerbatimInput[fontsize=\scriptsize,breaklines]{../.local/descifrado/transacciones.txt}%
}{%
 Este listado se incorpora al compilar con los resultados locales del script \texttt{scripts/descifrar.py}.%
}

\section*{Bibliografía}
\addcontentsline{toc}{section}{Bibliografía}
\href{https://drive.google.com/drive/folders/1rXMp5y9WeBR36iytc4RBaqXKWf3ea-WU}{Captura y claves de prueba del Laboratorio 4}.\par
\href{https://www.rfc-editor.org/rfc/rfc7914.html}{RFC 7914. The scrypt Password-Based Key Derivation Function}.\par
\medskip
\textbf{Repositorio Git:} \url{https://github.com/rkkas/laboratorio4-ssh}\par
\end{document}
